\documentclass[10pt,a4paper]{article}
\usepackage[margin=18mm]{geometry}
\usepackage[T1]{fontenc}
\usepackage{lmodern}
\usepackage{microtype}
\usepackage{amsmath,amssymb,mathtools,bm}
\usepackage{booktabs,longtable,tabularx,array}
\usepackage{siunitx}
\usepackage{graphicx,subcaption,float}
\usepackage{listings,fancyvrb}
\usepackage[dvipsnames]{xcolor}
\usepackage[colorlinks=true,allcolors=MidnightBlue]{hyperref}
\usepackage[nameinlink,noabbrev]{cleveref}
\usepackage{fancyhdr}
\usepackage{titlesec}
\usepackage{placeins}
\graphicspath{{../outputs/figures/}}
\sisetup{round-mode=places,round-precision=4,scientific-notation=true}
\allowdisplaybreaks
\pagestyle{fancy}
\setlength{\headheight}{14pt}
\fancyhf{}
\lhead{Wolfram Language derivation ledger}
\rhead{\thepage}
\newcommand{\RegressionPass}{42/42}
\newcommand{\RegressionMaxRel}{9.13e-07}
\newcommand{\RegressionMaxAbs}{2.43e-06}
\newcommand{\NMonteCarlo}{1200}
\newcommand{\DsLowG}{-0.061}
\newcommand{\DsLowWidth}{0.482}
\newcommand{\DsHighG}{-0.017}
\newcommand{\DsHighWidth}{0.094}
\newcommand{\BsLowG}{-0.548}
\newcommand{\BsLowWidth}{1.589}
\newcommand{\BsHighG}{-0.484}
\newcommand{\BsHighWidth}{2.332}


\definecolor{cellblue}{RGB}{232,240,248}
\definecolor{outputgray}{RGB}{245,245,245}
\newcommand{\InputCell}[1]{%
  \par\medskip\noindent\colorbox{cellblue}{%
  \parbox{\dimexpr\linewidth-2\fboxsep}{\texttt{In:}\quad #1}}\par\smallskip}
\newcommand{\OutputCell}[1]{%
  \par\smallskip\noindent\colorbox{outputgray}{%
  \parbox{\dimexpr\linewidth-2\fboxsep}{\texttt{Out:}\quad #1}}\par\medskip}
\lstdefinestyle{wl}{
  basicstyle=\ttfamily\scriptsize,
  keywordstyle=\color{MidnightBlue},
  commentstyle=\color{ForestGreen},
  stringstyle=\color{BrickRed},
  numbers=left,
  numberstyle=\tiny\color{gray},
  breaklines=true,
  breakatwhitespace=false,
  frame=single,
  rulecolor=\color{gray!45},
  columns=fullflexible,
  keepspaces=true,
  showstringspaces=false,
  tabsize=2
}
\lstset{style=wl}

\title{\textbf{Detailed Wolfram Language derivation and audit notebook}\\
Mixed-current $D_{s1},B_{s1}\to D_s^\ast,B_s^\ast\gamma$ sum rules}
\author{Machine-reproducible clean-room calculation}
\date{29 July 2026}

\begin{document}
\maketitle
\tableofcontents
\newpage

\section{Notebook status}

\InputCell{Evaluate the kernel, symbolic derivation, native numerical
regression, Python comparison, scans, and document generator.}
\begin{Verbatim}[fontsize=\small]
/Applications/Wolfram.app/Contents/MacOS/WolframKernel
  -noinit -noprompt -run 'Print[1+1];Quit[]'
Out: 2

Wolfram Kernel version: 15.0.0 for Mac OS X ARM (64-bit)
Random seed (Python uncertainty scan): 20260729
Regression: 42/42 entries pass
\end{Verbatim}
\OutputCell{The kernel is operational.  The symbolic and numerical Wolfram
Language paths complete without importing Python results.}

\begin{table}[H]
\centering\small
\begin{tabular}{lll}
\toprule
Layer & independent implementation & artifact\\
\midrule
symbolic Dirac/OPE & Wolfram Language & \path{outputs/analytic_formulas.wl}\\
symbolic report & Wolfram Language & \path{outputs/symbolic_validation.txt}\\
notebook & Wolfram Language & \path{notebooks/Dsstar_Bsstar_gamma2_derivation.nb}\\
numerical regression & Wolfram Language & \path{outputs/csv/mathematica_regression.csv}\\
full grids and MC & Python/NumPy & \path{outputs/csv/*.csv}\\
termwise comparison & Python/pandas & \path{outputs/csv/python_mathematica_regression.csv}\\
\bottomrule
\end{tabular}
\caption{Reproducibility layers.}
\end{table}

\section{Cell group 1: conventions and explicit gamma matrices}

\InputCell{Define the mostly-minus Dirac representation, $\gamma_5$,
$\sigma_{\mu\nu}$, and the Dirac adjoint.}
\begin{lstlisting}
z2 = ConstantArray[0, {2, 2}];
id2 = IdentityMatrix[2];
g0 = DiagonalMatrix[{1, 1, -1, -1}];
gSpatial[sig_] := ArrayFlatten[{{z2, sig}, {-sig, z2}}];
g = {g0, gSpatial[PauliMatrix[1]], gSpatial[PauliMatrix[2]],
     gSpatial[PauliMatrix[3]]};
g5 = I g[[1]].g[[2]].g[[3]].g[[4]];
sigma[mu_, nu_] :=
 I (g[[mu + 1]].g[[nu + 1]] - g[[nu + 1]].g[[mu + 1]])/2;
diracBar[x_] := g0.ConjugateTranspose[x].g0;
\end{lstlisting}
\OutputCell{$g_{\mu\nu}=(+,-,-,-)$,
$\gamma_5=i\gamma^0\gamma^1\gamma^2\gamma^3$,
$\epsilon^{0123}=+1$, and $p'=p+q$.}

The current basis and physical rotation are
\begin{align}
J_A^\mu&=\bar s\gamma^\mu\gamma_5Q,&
J_B^\mu&={i\over m_Q+m_s}\bar s\sigma^{\mu\nu}p'_\nu\gamma_5Q,\\
\begin{pmatrix}J_L\\J_H\end{pmatrix}
&=\begin{pmatrix}\sin\theta&\cos\theta\\
\cos\theta&-\sin\theta\end{pmatrix}
\begin{pmatrix}J_A\\J_B\end{pmatrix}.
\end{align}
The final current is $J_V^\mu=\bar s\gamma^\mu Q$.

\section{Cell group 2: cut traces and spectral matrix}

\InputCell{Place the two cut particles back-to-back in the $p'$ rest frame
and average the three transverse polarizations.}
\begin{lstlisting}
kSlash = g[[1]] eQ - g[[4]] kabs;
lSlash = g[[1]] es + g[[4]] kabs;
vertexA[i_] := g[[i + 1]].g5;
vertexB[i_] := I rootS sigma[i, 0].g5/(mQ + ms);
cutTrace[x_, y_, i_] :=
 Tr[(lSlash - ms IdentityMatrix[4]).x[i].
    (kSlash + mQ IdentityMatrix[4]).diracBar[y[i]]];
traceAAraw = FullSimplify[
 Sum[cutTrace[vertexA, vertexA, i], {i, 1, 3}]/3];
traceABraw = FullSimplify[
 Sum[cutTrace[vertexA, vertexB, i], {i, 1, 3}]/3];
traceBBraw = FullSimplify[
 Sum[cutTrace[vertexB, vertexB, i], {i, 1, 3}]/3];
\end{lstlisting}

Use
\begin{equation}
E_Q={s+m_Q^2-m_s^2\over2\sqrt s},\quad
E_s={s+m_s^2-m_Q^2\over2\sqrt s},\quad
|\bm k|^2={\lambda(s,m_Q^2,m_s^2)\over4s}.
\end{equation}
The phase-space normalization then yields
\begin{align}
\rho_{AA}&={\sqrt\lambda\over8\pi^2}
{[s-(m_Q+m_s)^2][(m_Q-m_s)^2+2s]\over s^2},\\
\rho_{AB}=\rho_{BA}&={\sqrt\lambda\over8\pi^2}
{3(m_Q-m_s)[s-(m_Q+m_s)^2]\over(m_Q+m_s)s},\\
\rho_{BB}&={\sqrt\lambda\over8\pi^2}
{[s-(m_Q+m_s)^2][s+2(m_Q-m_s)^2]\over(m_Q+m_s)^2s}.
\end{align}
\OutputCell{The AB and BA entries agree identically.  Every entry has mass
dimension two and the physical lower limit is $(m_Q+m_s)^2$.}

\InputCell{Factor the determinant of the spectral matrix.}
\begin{lstlisting}
spectralDeterminant =
 Factor[rhoAA rhoBB - rhoAB^2]
\end{lstlisting}
\OutputCell{\[
\det\rho={\left[s-(m_Q+m_s)^2\right]^3
\left[s-(m_Q-m_s)^2\right]^3
\over32\pi^4(m_Q+m_s)^2s^3}\ge0 .
\]}

\section{Cell group 3: Borel matrix and physical residues}

\InputCell{Construct the continuum-subtracted perturbative matrix and add
the explicit dimension-three, finite-$m_s$, and dimension-five matrices.}
\begin{align}
\widehat\Pi^{(3)}&=\langle\bar ss\rangle e^{-m_Q^2/M^2}
\begin{pmatrix}
m_Q&m_Q^2/S\\m_Q^2/S&m_Q^3/S^2
\end{pmatrix},\\
\widehat\Pi^{(m_s)}&=\langle\bar ss\rangle e^{-m_Q^2/M^2}
\left(-{m_Q^2m_s\over2M^2}+{m_Q^3m_s^2\over2M^4}\right)
\begin{pmatrix}1&m_Q/S\\m_Q/S&m_Q^2/S^2\end{pmatrix},\\
\widehat\Pi^{(5)}&=e^{-m_Q^2/M^2}
\begin{pmatrix}
-m_0^2\langle\bar ss\rangle m_Q^3/(4M^4)&
-{m_0^2\langle\bar ss\rangle\over4S}
(m_Q^4/M^4-m_Q^2/M^2-1)\\
\mathrm{sym.}&
{m_0^2\langle\bar ss\rangle\over S^2}
[-m_Q^5/(4M^4)+m_Q^3/(2M^2)]
\end{pmatrix}.
\end{align}
Here $S=m_Q+m_s$.  The final basis matrix is
\begin{equation}
\widehat\Pi(M^2,s_0)=
\int_{S^2}^{s_0}ds\,e^{-s/M^2}
\begin{pmatrix}\rho_{AA}&\rho_{AB}\\\rho_{AB}&\rho_{BB}\end{pmatrix}
+\widehat\Pi^{(3)}+\widehat\Pi^{(m_s)}+\widehat\Pi^{(5)}.
\end{equation}

\InputCell{Rotate the complete matrix; do not diagonalize or zero AB by hand.}
\begin{lstlisting}
rotation = {{Sin[theta], Cos[theta]},
            {Cos[theta], -Sin[theta]}};
piPhysical = Expand[
  rotation.piBasis.Transpose[rotation]]
\end{lstlisting}
\OutputCell{\[
\Pi_{LL}=s_\theta^2\Pi_{AA}+2s_\theta c_\theta\Pi_{AB}
+c_\theta^2\Pi_{BB},\quad
\Pi_{HH}=c_\theta^2\Pi_{AA}-2s_\theta c_\theta\Pi_{AB}
+s_\theta^2\Pi_{BB}.
\]}

The residue and derivative mass checks are
\begin{equation}
f_a={e^{m_a^2/(2M^2)}\sqrt{\Pi_{aa}^{\rm phys}}\over m_a},
\qquad
m_{\rm SR}^2={-\partial_{1/M^2}\Pi_{aa}^{\rm phys}
\over\Pi_{aa}^{\rm phys}}.
\end{equation}
The accepted windows are reproduced directly from the full scan:
\begin{table}[H]
\centering\small
\begin{tabular}{lcccccc}
\toprule
State & $M^2$ [GeV$^2$] & $s_0$ [GeV$^2$] & accepted & PC & $|d=5|/|\Pi|$ & $m_{\rm SR}$ [GeV]\\
\midrule
$D_{s1}(2460)$ & 2.0--3.0 & 7.50--8.50 & 29/30 & 0.41--0.71 & 0.023 & 2.418--2.506 \\
$D_{s1}(2536)$ & 2.0--3.0 & 8.50--10.00 & 18/30 & 0.41--0.67 & 0.097 & 2.428--2.655 \\
$B_{s1}^{L}$ & 8.0--11.0 & 40.00--42.00 & 30/30 & 0.42--0.66 & 0.002 & 5.649--5.785 \\
$B_{s1}(5830)$ & 7.0--9.0 & 42.00--44.00 & 30/30 & 0.40--0.64 & 0.037 & 5.730--5.896 \\
\bottomrule
\end{tabular}

\caption{Two-point acceptance ledger.}
\end{table}

\section{Cell group 4: vector final state}

\InputCell{Repeat the sum-rule extraction for
$J_V^\mu=\bar s\gamma^\mu Q$.}
\begin{equation}
\rho_V(s)={\sqrt\lambda\over8\pi^2}
\left[2-{m_Q^2+m_s^2-6m_Qm_s\over s}
-{(m_Q^2-m_s^2)^2\over s^2}\right].
\end{equation}
The local AA contribution changes sign under the parity flip.  The same pole
and local-fraction gates are applied.  The Monte Carlo vector residues are
listed together with the initial residues in \cref{tab:nbresults}.

\section{Cell group 5: external-photon LCSR}

\InputCell{Project the correlator onto
$ie\,\epsilon_{\alpha\beta\rho\sigma}
\eta^\alpha\xi^{*\beta}\varepsilon^{*\rho}q^\sigma$.}
\begin{equation}
T_{\mu\nu}^i=i\int d^4x\,e^{ipx}
\langle\gamma|T\{J_{V\mu}^\dagger(x)J_\nu^i(0)\}|0\rangle .
\end{equation}
The Wick contraction supplies the global minus sign:
\begin{equation}
T_{\mu\nu}^i=-i\int d^4x\,e^{ipx}
\langle\gamma|\bar s(x)\gamma_\mu S_Q(x,0)\Gamma_\nu^i s(0)|0\rangle.
\end{equation}
Hard emission from both charged lines and nonlocal photon DAs through twist
four are included.  No ordinary local transition condensate is included.

\InputCell{Apply the equal double-Borel transform.}
\begin{equation}
{\cal B}_{M_1^2}{\cal B}_{M_2^2}{1\over
[m_Q^2-(1-u)p^2-up'^2]^n}
={({\cal M}^2)^{2-n}\over\Gamma(n)}
e^{-m_Q^2/{\cal M}^2}\delta(u-u_0),
\end{equation}
where $u_0=M_1^2/(M_1^2+M_2^2)$ and
${\cal M}^2=M_1^2M_2^2/(M_1^2+M_2^2)$.  The choice
$M_1^2=M_2^2=2M^2$ gives $u_0=1/2$ and ${\cal M}^2=M^2$.

With $E=e^{-m_Q^2/M^2}$ and $E_0=e^{-s_0/M^2}$, the six included terms are
\begin{align}
\widehat T_{\rm hard}&=\int_{(m_Q+m_s)^2}^{s_0}ds\,
e^{-s/M^2}\rho_T(s),\\
\widehat T_{\rm tw2}&=e_sm_Q\langle\bar ss\rangle\chi
M^2(E-E_0)\phi_\gamma(u_0),\\
\widehat T_{\rm tw4,2p}&=e_sm_Q\langle\bar ss\rangle E
\left[-{m_Q^2{\cal A}\over4M^2}-(1-u_0)H_\gamma-\bar H_\gamma
+{2m_Q^2\bar H_\gamma\over M^2}\right],\\
\widehat T_{\rm tw3,2p}&=e_sf_{3\gamma}M^2(E-E_0)
\left[{(1-u_0)\psi^{a\prime}\over4}-{\psi^a\over4}
-(1+2m_Q^2/M^2)\Psi^v+(1-u_0)\psi^v\right],\\
\widehat T_{\rm tw4,3p}&=m_Qe_s\langle\bar ss\rangle E I_{{\cal F}_2},\\
\widehat T_{\rm tw3,3p}&=-e_sf_{3\gamma}M^2(E-E_0) I_{{\cal F}_3}.
\end{align}
The hard density is
\begin{equation}
\rho_T={3m_sm_Q\over4\pi^2}\left[
e_s\ln{a_s-\sqrt\lambda\over a_s+\sqrt\lambda}
+e_Q\ln{a_Q-\sqrt\lambda\over a_Q+\sqrt\lambda}\right],
\end{equation}
$a_s=s-m_Q^2+m_s^2$ and $a_Q=s-m_s^2+m_Q^2$.

\begin{table}[H]
\centering\scriptsize
\resizebox{\textwidth}{!}{\begin{tabular}{lrrrrrrr}
\toprule
State & hard & tw2 & tw3(2p) & tw3(3p) & tw4(2p) & tw4(3p) & local diag.\\
\midrule
$D_{s1}(2460)$ & \num{0.019} & \num{-0.051} & \num{0.013} & \num{0.013} & \num{-8.631e-04} & \num{-6.133e-21} & \num{-0.007} \\
$D_{s1}(2536)$ & \num{0.020} & \num{-0.054} & \num{0.014} & \num{0.014} & \num{-8.631e-04} & \num{-6.133e-21} & \num{-0.007} \\
$B_{s1}^{L}$ & \num{0.013} & \num{-0.023} & \num{0.002} & \num{0.002} & \num{-0.004} & \num{-1.515e-21} & \num{8.400e-04} \\
$B_{s1}(5830)$ & \num{0.013} & \num{-0.023} & \num{0.002} & \num{0.002} & \num{-0.004} & \num{-1.515e-21} & \num{8.400e-04} \\
\bottomrule
\end{tabular}
}
\caption{Numerical transition OPE ledger in $\mathrm{GeV}^4$.  The local
diagnostic is excluded.}
\end{table}

\section{Cell group 6: tensor-current projection}

\InputCell{Evaluate the twist-two Dirac traces for A and B before Borel
transformation.}
\begin{lstlisting}
twist2TraceA = -8 mQ;
twist2TraceB = -8 k2/(mQ + ms);
twist2PreBorelRatio =
 Factor[twist2TraceB/twist2TraceA]
\end{lstlisting}
\OutputCell{\[
{T_B^{(2)}\over T_A^{(2)}}={k^2\over m_Q(m_Q+m_s)}.
\]}
The identity
$k^2/(m_Q^2-k^2)=m_Q^2/(m_Q^2-k^2)-1$ separates a contact term whose Borel
image is zero.  Hence
\begin{equation}
\widehat T_B={m_Q\over m_Q+m_s}\widehat T_A.
\end{equation}
This leading-power projection is used term by term.  The uncalculated
subleading tensor kernels carry an explicit 15\% charm or 5\% bottom
correlated nuisance.

\section{Cell group 7: physical coupling and width}

For row $a$ of the rotation matrix,
\begin{equation}
\widehat T_a=R_{aA}\widehat T_A+R_{aB}\widehat T_B,\qquad
g_a={e^{(m_a^2+m_V^2)/(2M^2)}\widehat T_a
\over f_af_Vm_am_V}.
\end{equation}
The amplitude and width are
\begin{align}
{\cal A}&=ie\,g_a\,\epsilon_{\alpha\beta\rho\sigma}
\eta^\alpha\xi^{*\beta}\varepsilon^{*\rho}q^\sigma,\\
\Gamma&={\alpha_{\rm em}g_a^2k_\gamma^3(m_a^2+m_V^2)
\over3m_a^2m_V^2}.
\end{align}
\InputCell{Replace $\varepsilon\to q$ in the determinant form of the
Levi-Civita contraction.}
\begin{lstlisting}
wardDeterminant =
 Det[{{eta0, eta1, eta2, eta3},
      {xi0, xi1, xi2, xi3},
      {q0, q1, q2, q3},
      {q0, q1, q2, q3}}];
FullSimplify[wardDeterminant] === 0
\end{lstlisting}
\OutputCell{True.  The retained structure satisfies the Ward identity
exactly.}

\section{Cell group 8: numerical results and non-Gaussian uncertainty}

\begin{table}[H]
\centering\small
\resizebox{\textwidth}{!}{\begin{tabular}{lcccc}
\toprule
State & $g$ & $\Gamma_\gamma$ [keV] & $f_i$ [GeV] & $f_V$ [GeV]\\
\midrule
$D_{s1}(2460)$ & $-0.061_{-0.160}^{+0.174}$ & $0.482_{-0.434}^{+1.609}$ & $0.404_{-0.018}^{+0.020}$ & $0.253_{-0.010}^{+0.008}$ \\
$D_{s1}(2536)$ & $-0.017_{-0.069}^{+0.058}$ & $0.094_{-0.089}^{+0.520}$ & $0.166_{-0.012}^{+0.014}$ & $0.253_{-0.010}^{+0.008}$ \\
$B_{s1}^{L}$ & $-0.548_{-0.339}^{+0.320}$ & $1.589_{-1.298}^{+2.703}$ & $0.501_{-0.017}^{+0.018}$ & $0.300_{-0.009}^{+0.009}$ \\
$B_{s1}(5830)$ & $-0.484_{-0.425}^{+0.286}$ & $2.332_{-1.925}^{+5.735}$ & $0.098_{-0.008}^{+0.014}$ & $0.300_{-0.009}^{+0.009}$ \\
\bottomrule
\end{tabular}
}
\caption{Medians and 16th--84th percentiles from 1200 accepted samples per
state.}
\label{tab:nbresults}
\end{table}

The median widths are \DsLowWidth, \DsHighWidth, \BsLowWidth, and
\BsHighWidth\ keV in the state order of \cref{tab:nbresults}.  The charm
coupling intervals cross zero, and both higher states show destructive
current interference.  The width distributions are therefore non-Gaussian.

\begin{table}[H]
\centering\small
\begin{tabular}{lrrrr}
\toprule
State & $\Delta_{M^2}g$ & $\Delta_{s_0}g$ & $\Delta_{M^2}\Gamma$ & $\Delta_{s_0}\Gamma$\\
\midrule
$D_{s1}(2460)$ & 39.8\% & 3.4\% & 87.1\% & 6.7\% \\
$D_{s1}(2536)$ & 40.7\% & 3.3\% & 89.3\% & 6.6\% \\
$B_{s1}^{L}$ & 108.6\% & 1.2\% & 279.9\% & 2.3\% \\
$B_{s1}(5830)$ & 113.4\% & 0.9\% & 295.7\% & 1.9\% \\
\bottomrule
\end{tabular}

\caption{Full-range relative variation $(\max-\min)/|\mathrm{mean}|$.}
\end{table}
\OutputCell{The $s_0$ scans are mild; the transition $M^2$ scans are not.
The bottom widths vary by roughly 280--296\% across the declared window.
These are exploratory predictions, not precision plateaus.}

\section{Cell group 9: independent numerical regression}

\InputCell{Evaluate the same selected densities, Borel matrices, transition
terms, residues, couplings, and widths with native Mathematica
\texttt{NIntegrate}.  Compare with the Python Gauss--Legendre implementation.}
\begin{longtable}{llrrr}
\toprule
sector & quantity & Python & Mathematica & relative difference\\
\midrule
\endfirsthead
\toprule
sector & quantity & Python & Mathematica & relative difference\\
\midrule
\endhead
D & rho\_AA & \num{0.013362} & \num{0.013362} & 0\\
D & rho\_AB & \num{0.013932} & \num{0.013932} & 0\\
D & rho\_BB & \num{0.016293} & \num{0.016293} & 0\\
D & T\_quoted\_total\_A & \num{-0.006170} & \num{-0.006170} & \num{4.39e-07}\\
D & T\_quoted\_total\_B & \num{-0.005749} & \num{-0.005749} & \num{4.39e-07}\\
D & Ds1\_2460\_g & \num{-0.044925} & \num{-0.044925} & \num{4.56e-07}\\
D & Ds1\_2460\_width\_keV & \num{0.064298} & \num{0.064298} & \num{9.13e-07}\\
D & Ds1\_2536\_g & \num{-0.019738} & \num{-0.019738} & \num{4.50e-07}\\
D & Ds1\_2536\_width\_keV & \num{0.020963} & \num{0.020963} & \num{9.00e-07}\\
B & rho\_AA & \num{0.003014} & \num{0.003014} & \num{3.31e-14}\\
B & rho\_AB & \num{0.003016} & \num{0.003016} & 0\\
B & rho\_BB & \num{0.003032} & \num{0.003032} & \num{3.29e-14}\\
B & T\_quoted\_total\_A & \num{-0.010237} & \num{-0.010237} & \num{3.73e-07}\\
B & T\_quoted\_total\_B & \num{-0.010015} & \num{-0.010015} & \num{3.73e-07}\\
B & Bs1\_lower\_g & \num{-0.531512} & \num{-0.531512} & \num{3.87e-07}\\
B & Bs1\_lower\_width\_keV & \num{1.515962} & \num{1.515963} & \num{7.74e-07}\\
B & Bs1\_5830\_g & \num{-0.565445} & \num{-0.565445} & \num{3.87e-07}\\
B & Bs1\_5830\_width\_keV & \num{3.130978} & \num{3.130980} & \num{7.75e-07}\\
\bottomrule
\end{longtable}

\OutputCell{\RegressionPass\ entries pass.  Maximum material relative
difference: \RegressionMaxRel.  Maximum absolute difference:
\RegressionMaxAbs.}

\section{Cell group 10: validation gate ledger}

\begin{longtable}{p{0.29\linewidth}p{0.13\linewidth}p{0.50\linewidth}}
\toprule
Gate & status & evidence\\
\midrule
\endfirsthead
\toprule
Gate & status & evidence\\
\midrule
\endhead
$\rho_{AB}=\rho_{BA}$ & checked & identical function and symbolic WL equality\\
spectral positivity & checked & analytic determinant is nonnegative above threshold\\
mass dimensions & checked & rho dim 2; Borel matrix and T dim 4; g dimensionless\\
physical threshold & checked & all dispersive integrals start at (mQ+ms)\textasciicircum{}2\\
Ward identity & checked & epsilon structure vanishes under epsilon\_gamma -$>$ q\\
no local transition condensate & checked & diagnostic term excluded from quoted\_total\_A/B\\
two-point $d=3,5$ & checked & explicit matrices retained\\
pure A limit & checked & theta=90 degrees for lower state\\
pure B limit & checked & theta=0 degrees for lower state\\
equal Borel limit & checked & M1\textasciicircum{}2=M2\textasciicircum{}2=2M\textasciicircum{}2 gives u0=1/2\\
$m_s\to0$ threshold & checked & lambda and density reduce smoothly\\
charge switch & checked & hard density changes through eQ/es only\\
Python--Mathematica regression & checked & 42/42 entries pass; max material rel=9.128e-07; max abs=2.425e-06\\
subleading tensor kernels & truncation & not retained; 15\% charm and 5\% bottom correlated uncertainty assigned\\
transition $\alpha_s$ & truncation & not retained\\
\bottomrule
\end{longtable}


\section{Complete figure atlas}

Every plotted PDF below has a PNG twin and a companion CSV.  The sequence is
fixed across states: $D_{s1}(2460)$, $D_{s1}(2536)$, $B_{s1}^{L}$, and
$B_{s1}(5830)$.

\newcommand{\FourPlots}[5]{%
\begin{figure}[p]\centering
\begin{subfigure}{0.48\textwidth}\includegraphics[width=\linewidth]{Ds1_2460_#1.pdf}\caption{$D_{s1}(2460)$}\end{subfigure}
\begin{subfigure}{0.48\textwidth}\includegraphics[width=\linewidth]{Ds1_2536_#1.pdf}\caption{$D_{s1}(2536)$}\end{subfigure}\\
\begin{subfigure}{0.48\textwidth}\includegraphics[width=\linewidth]{Bs1_lower_#1.pdf}\caption{$B_{s1}^{L}$}\end{subfigure}
\begin{subfigure}{0.48\textwidth}\includegraphics[width=\linewidth]{Bs1_5830_#1.pdf}\caption{$B_{s1}(5830)$}\end{subfigure}
\caption{#2.  #3}\label{#4}
\end{figure}\FloatBarrier}

\FourPlots{coupling_vs_M2}{Coupling versus transition Borel mass}{$s_0$ fixed
at its central value.}{fig:all-gm2}{}
\FourPlots{width_vs_M2}{Width versus transition Borel mass}{$s_0$ fixed at
its central value.}{fig:all-wm2}{}
\FourPlots{coupling_vs_s0}{Coupling versus continuum threshold}{$M^2$ fixed
at its central value.}{fig:all-gs0}{}
\FourPlots{width_vs_s0}{Width versus continuum threshold}{$M^2$ fixed at its
central value.}{fig:all-ws0}{}
\FourPlots{coupling_vs_angle}{Coupling versus mixing angle}{The selected
angle and pure-current endpoints are displayed.}{fig:all-angle}{}
\FourPlots{pc_vs_M2}{Two-point pole contribution}{Central two-point threshold.}
{fig:all-pc}{}
\FourPlots{mass_sr_vs_M2}{Two-point reconstructed mass}{Horizontal reference
is the physical or adopted lower-state mass.}{fig:all-mass}{}
\FourPlots{f_vs_M2}{Two-point current residue}{Central two-point threshold.}
{fig:all-f}{}
\FourPlots{d5frac_vs_M2}{Dimension-five OPE fraction}{Central two-point
threshold.}{fig:all-d5}{}

\section{Raw symbolic output}

\InputCell{Import the text artifact written directly by the symbolic
Wolfram Language script.}
\lstinputlisting{../outputs/symbolic_validation.txt}

\appendix
\section{Complete symbolic Wolfram Language source}

The listing below is the source that generated the notebook, analytic formula
association, and symbolic validation report.
\lstinputlisting{../scripts/symbolic_derivation.wl}

\section{Complete independent Mathematica regression source}

This listing evaluates native numerical integrals and writes the Mathematica
CSV.  It never reads a Python result.
\lstinputlisting{../scripts/mathematica_regression.wl}

\section{Reproduction commands}

\begin{lstlisting}
export PROJECT_PYTHON=/Users/sbilmis/.cache/codex-runtimes/codex-primary-runtime/dependencies/python/bin/python3
"$PROJECT_PYTHON" scripts/run_analysis.py
/Applications/Wolfram.app/Contents/MacOS/WolframKernel -noinit -noprompt -script scripts/symbolic_derivation.wl
/Applications/Wolfram.app/Contents/MacOS/WolframKernel -noinit -noprompt -script scripts/mathematica_regression.wl
"$PROJECT_PYTHON" scripts/compare_regression.py
"$PROJECT_PYTHON" scripts/build_documents.py
latexmk -pdf -interaction=nonstopmode -halt-on-error -cd tex/Dsstar_Bsstar_gamma2_paper.tex
latexmk -pdf -interaction=nonstopmode -halt-on-error -cd tex/Dsstar_Bsstar_gamma2_Mathematica_derivation.tex
\end{lstlisting}

\end{document}
